\section{Electric-flux localization and restricted pair denominators}
\label{flux4-section}

We retain the compact-link Hamiltonian~\eqref{lat-hamiltonian},
with no electric truncation in its definition.  In this section
the periodic cubic lattice has $n\ge3$, $N=n^3$ sites and $3N$ positive
links.  Put
\begin{equation}
 c=\frac{g^2}{2a},\qquad b=\frac1{g^2a},\qquad
 \mu=m+\frac{3r}{a},\qquad
 h_*=m+\frac{3+6r}{a}.
 \label{flux4-constants}
\end{equation}
All assertions below concern the stated finite lattice, with $g^2>0$,
$a>0$, $m,r\ge0$.  Positive lower bounds on pair denominators will be
separate hypotheses.  A function of the commuting integer electric
operators commutes with every Gauss operator
$G_x=\operatorname{div}E_x+q_x$.  Consequently all electric cutoffs used
below preserve each physical Gauss sector.  This does not make their
ranges invariant under $H$.

\subsection{The actual vacuum has electric tails}

The Fourier representation of the gauge Hilbert space is
$\ell^2(\mathbb Z^{\mathcal E})$.  If $s\in\mathbb Z^{\mathcal E}$,
write $U^s=\prod_\ell U_\ell^{s_\ell}$, so
$U^s|e\rangle=|e+s\rangle$.  Each fermion hopping changes one electric
coordinate by one.  Each plaquette multiplier changes the four
coordinates on its oriented boundary by one, with signs.  Denote this
plaquette vector by $s_p$; it has $|s_p|^2=4$ and
$\operatorname{div}s_p=0$.

\begin{proposition}[No eigenvector has finite electric support]
\label{flux4-vacuum-tails}
Every nonzero eigenvector of $H$, including an eigenvector in an
admissible fixed Gauss sector, has infinite support in the electric
Fourier variables.  In particular every normalized physical vacuum
ground vector $\Omega$ satisfies, for every finite $F>0$,
\begin{equation}
 \left\langle\Omega,
   \sum_{\ell\in\mathcal E}{\bf1}_{\{|E_\ell|\ge F\}}
   \Omega\right\rangle>0.
 \label{flux4-strict-tail}
\end{equation}
\end{proposition}
\begin{proof}
The nonconstant part of the magnetic multiplier is
$-\frac b2\sum_p(U^{s_p}+U^{-s_p})$.  Choose a plaquette vector $s_p$.
The linear functional $e\mapsto s_p\cdot e$ takes its strictly largest
value on the set of distinct nonzero shifts of $H$ at $s_p$: its value
there is $4$, its value on a single-link shift is at most $1$, and its
value on another plaquette vector is less than $4$ unless that vector
equals $s_p$.  This follows from Cauchy--Schwarz and equality in it.
If a plaquette vector occurs more than once, its scalar coefficients
add with the same negative sign.

Suppose an eigenvector $\psi$ had finite electric support $S$.
Its Fourier coefficient $\psi_e$ is a vector in the finite fermion
Fock space.  Perturb the above linear functional slightly so that
it still uniquely maximizes the distinct shifts at $s_p$ and uniquely
maximizes $S$ at an element $e_*$.  Such a perturbation exists because
only finitely many strict inequalities and hyperplanes are involved.
The coefficient of $H\psi$ at $e_*+s_p$ then has exactly one possible
shift contribution: a nonzero scalar multiple of $\psi_{e_*}$ from
the plaquette multiplier.  Indeed, any contributing $e+s=e_*+s_p$
would have to maximize separately both linear functionals on $S$
and on the shifts.  The electric kinetic term and all zero-shift
terms contribute zero, since $e_*+s_p\notin S$.  Thus this coefficient
is nonzero, whereas the same coefficient of an eigenvalue times
$\psi$ is zero.  This is a contradiction.

If the nonnegative expectation in \eqref{flux4-strict-tail} vanished,
every Fourier coefficient outside a finite electric cube would vanish.
The first assertion therefore proves the second.  This argument uses
the actual magnetic plaquette multiplier and the full fermionic
Hamiltonian, rather than a scalar rotor surrogate.
\end{proof}

\begin{corollary}[Failure of a source-only flux budget]
\label{flux4-failed-relative-budget}
No finite constant can make the implication
\begin{equation}
 \langle H\rangle\le E_0+C(Z_{\rm tot}+K)\quad\Longrightarrow\quad
 \left\langle\sum_\ell{\bf1}_{\{|E_\ell|\ge F\}}\right\rangle
 \le\frac{C'(Z_{\rm tot}+K)}{cF^2}
 \label{flux4-false-budget}
\end{equation}
hold for all normalized physical states and finite $F$.  It fails
already at $K=Z_{\rm tot}=0$ for a vacuum ground vector.
\end{corollary}
\begin{proof}
The energy hypothesis is then an equality, the proposed right side
is zero, and \eqref{flux4-strict-tail} makes the left side positive.
\end{proof}

Subtracting the scalar vacuum energy from $H$ does not subtract the
vacuum fluctuations of $E_\ell^2$.  Nor does positivity of $H_E$ give
positivity of $\langle H_E\rangle_\psi-\langle H_E\rangle_\Omega$.
A theorem controlling an excess flux distribution would need a
separate response estimate and an explicitly chosen reference state.

\begin{lemma}[A valid extensive flux budget]
\label{flux4-volume-budget}
For any normalized form-domain vector with
$\langle H\rangle\le E_0+\Delta$, one has
\begin{align}
 cF^2\left\langle\sum_\ell
       {\bf1}_{\{|E_\ell|\ge F\}}\right\rangle
 &\le \langle H_E\rangle
 \le E_0+\Delta+2Nh_* \notag\\
 &\le \Delta+N\left(\frac{6(1+r)}a+\frac3{g^2a}\right).
 \label{flux4-volume-budget-eq}
\end{align}
These inequalities apply in every physical sector.  In particular,
for the sharp cutoff
$P_F=\prod_\ell{\bf1}_{\{|E_\ell|\le F\}}$ with integer $F\ge0$,
\begin{equation}
 \|(1-P_F)\psi\|^2
 \le\frac{\Delta+N[6(1+r)/a+3/(g^2a)]}{c(F+1)^2}.
 \label{flux4-outside-probability}
\end{equation}
\end{lemma}
\begin{proof}
The first inequality is pointwise on the electric joint spectrum.
The bound $d\Gamma(h_W(U))\ge-2Nh_*$ and $H_B\ge0$ give
$H\ge H_E-2Nh_*$.  Finally use the vacuum trial bound
$E_0\le-2N\mu+3N/(g^2a)$ from Lemma~\ref{lat-realization}.
The last displayed estimate follows because outside $P_F$ at least
one integer electric coordinate has magnitude at least $F+1$.
\end{proof}

This is a uniform bound per site, not a bound proportional only to
the sources.  It remains valid when $\Delta=C(Z_{\rm tot}+K)$, but
retains its volume term.  The bare-excitation obstruction in
the result~\ref{rel4-bare-sea} is the fermionic counterpart: the
interacting vacuum is not annihilated by the onsite excitation number.
Neither obstruction excludes a suitably dressed, source-relative
estimate; both exclude the proposed estimates for the bare observables.

\paragraph{Gauss's law supplies no upper flux bound near a nucleus.}
At a site with source $\nu_x\ge2$, the six signed incident electric
coordinates have sum $\nu_x-q_x$, where $q_x\le2$.  Thus on that
physical sector
\begin{equation}
 \sum_{\ell\ni x}E_\ell^2\ge\frac{(\nu_x-2)^2}{6}.
 \label{flux4-local-gauss-lower}
\end{equation}
For $\nu_x=26$ this is $96$, and at least one incident coordinate
has magnitude at least $4$ in each electric--occupation basis vector.
These are lower bounds.  Adding an arbitrarily large integer
divergence-free plaquette circulation preserves Gauss's law and
precludes a kinematical upper bound of order $Z/6$.  An upper bound
in low-energy states must be proved using their energy or distribution.

\subsection{A Gauss-preserving electric IMS formula}

For explicit constants write
\begin{equation}
 H=H_E+C_0+\sum_{s\in\mathcal S}U^s A_s,
 \qquad
 \mathcal J:=\sum_{s\in\mathcal S}\|A_s\|
 \le\frac{12N(1+r)}a+\frac{3N}{g^2a}.
 \label{flux4-shift-decomposition}
\end{equation}
Here a repeated shift may be listed repeatedly, $C_0$ acts only on
fermions and is independent of the electric variables, and every
shift has $\|s\|_\infty=1$.  The $A_s$ are finite fermion matrices
(scalar matrices for the plaquettes).  To verify the stated bound,
the directed fermion coefficient on a link is
\[
 A_\ell=\sum_{u,v=1}^4c_{xu}^*
       \left(\frac{-i\alpha_j-r\beta}{2a}\right)_{uv}c_{yv}.
\]
Singular-value decomposition and the CAR give
$\|A_\ell\|\le\|(-i\alpha_j-r\beta)/(2a)\|_1
\le2(1+r)/a$.  There are two such shifts per positive link.
The two magnetic coefficients of a plaquette each have norm $b/2$.
This proves \eqref{flux4-shift-decomposition}; no commutativity of
different fermion coefficients is used.

Choose a Lipschitz function $\theta:[0,\infty)\to[0,\pi/2]$ which
is zero on $[0,1]$ and equals $\pi/2$ on $[2,\infty)$, and write
$L_\theta=\operatorname{Lip}(\theta)$.  It can equally well be chosen
smooth.  Let $R(E)=\max_\ell|E_\ell|$, and set
\begin{equation}
 f_0(E)=\cos\theta(R(E)/F),\qquad
 f_1(E)=\sin\theta(R(E)/F),\qquad F\ge1.
 \label{flux4-two-cutoffs}
\end{equation}
Then $f_0^2+f_1^2=1$, the range of $f_0$ lies in the electric cube
$R\le2F$, and $f_1$ vanishes on $R\le F$.

\begin{lemma}[Electric IMS identity with a quadratic cutoff error]
\label{flux4-ims}
The $f_j$ preserve $D(H_E)$, its form domain, and each Gauss sector.
There is a bounded self-adjoint operator $\mathcal R_F$ such that
\begin{equation}
 f_0Hf_0+f_1Hf_1=H+\mathcal R_F,
 \qquad
 \|\mathcal R_F\|\le
       \frac{L_\theta^2\mathcal J}{2F^2}.
 \label{flux4-ims-bound}
\end{equation}
The identity holds both on $D(H_E)$ and as a quadratic-form identity.
With
$B_F={\bf1}_{\{F-1\le R(E)\le2F+1\}}$, one also has
$\mathcal R_F=B_F\mathcal R_FB_F$ and therefore
\begin{equation}
 |\langle\psi,\mathcal R_F\psi\rangle|
 \le\frac{L_\theta^2\mathcal J}{2F^2}
       \langle\psi,B_F\psi\rangle.
 \label{flux4-ims-shell}
\end{equation}
\end{lemma}
\begin{proof}
The cutoffs commute with $H_E$, $C_0$, and every $G_x$.
For one shift,
\[
 \sum_{j=0}^1 f_jU^sA_sf_j-U^sA_s
 =U^sA_s\left[
   \cos\big(\theta(R(E+s)/F)-\theta(R(E)/F)\big)-1\right].
\]
Since $|R(e+s)-R(e)|\le1$ and $1-\cos t\le t^2/2$,
the bracket has norm at most $L_\theta^2/(2F^2)$.
Summing the finitely many shifts proves the norm bound.
The bracket can be nonzero only when both endpoints lie in
$F-1\le R\le2F+1$, which proves the shell assertion.
All unbounded terms commute with the cutoffs, so the core calculation
extends directly to the stated domains.  Self-adjointness of the
bounded error also follows from the left side of
\eqref{flux4-ims-bound}.
\end{proof}

There is a useful fully local variant.  For each link set
$u_{0\ell}=\cos\theta(|E_\ell|/F)$ and
$u_{1\ell}=\sin\theta(|E_\ell|/F)$.  The $2^{3N}$ products
$f_\sigma=\prod_\ell u_{\sigma_\ell,\ell}$ form a quadratic
partition of unity and commute with all Gauss operators.  Define
$B_{\ell,F}={\bf1}_{\{F-1\le|E_\ell|\le2F+1\}}$ and
\begin{equation}
 j_{\rm loc}=\frac{4(1+r)}a+\frac4{g^2a}.
 \label{flux4-local-shift-constant}
\end{equation}
Each link meets four plaquettes, so
$\sum_{s:\,s_\ell\ne0}\|A_s\|\le j_{\rm loc}$.

\begin{lemma}[Local electric IMS error]
\label{flux4-local-ims}
The product partition satisfies
\begin{equation}
 \left|\left\langle\psi,
 \left(\sum_\sigma f_\sigma Hf_\sigma-H\right)\psi\right\rangle\right|
 \le\frac{L_\theta^2j_{\rm loc}}{2F^2}
       \sum_\ell\langle\psi,B_{\ell,F}\psi\rangle.
 \label{flux4-local-ims-bound}
\end{equation}
\end{lemma}
\begin{proof}
The overlap of the two product vectors at electric configurations
$e,e+s$ is $\prod_{\ell:s_\ell\ne0}\cos d_\ell$, where
$d_\ell=\theta(|e_\ell+s_\ell|/F)-\theta(|e_\ell|/F)$.
These cosines lie in $[0,1]$, so
\[
 0\le1-\prod_\ell\cos d_\ell
 \le\sum_\ell(1-\cos d_\ell)
 \le\frac{L_\theta^2}{2F^2}
       \sum_{\ell:s_\ell\ne0}
       B_{\ell,F}(e+s)B_{\ell,F}(e).
\]
For each summand of the Fourier kernel, bound the fermion matrix by
$\|A_s\|$ and apply Cauchy--Schwarz to the sum in $e$.  The term
with the two boundary indicators is at most
$\|A_s\|\langle\psi,B_{\ell,F}\psi\rangle$ times
$L_\theta^2/(2F^2)$.  Summing first over shifts touching a fixed link
gives \eqref{flux4-local-ims-bound}.
\end{proof}

These identities are valid localization statements, including in
nonzero source sectors.  The local partition does not discard its
branches containing high-flux links.  The global bound
\eqref{flux4-outside-probability} makes the discarded norm small by
this argument only when $F^2$ is large compared with $N$ at fixed
parameters.  A positive pair denominator, in contrast, puts an upper
bound on $F$ at fixed parameters.  Thus these bounds alone do not
produce a uniform macroscopic low-energy reduction.

\subsection{Sharp truncations and their boundary terms}

For integer $F\ge0$, compress $H$ to $\operatorname{ran}P_F$ and
then to the chosen Gauss sector.  Denote this finite matrix by $H_F$;
an empty compressed physical sector is excluded.  The finite matrices
are mathematically well defined, but need not have the positive
electric heat kernel used to prove Theorem~\ref{lat-main}.

\begin{lemma}[Convergence and exact sharp-boundary coupling]
\label{flux4-sharp-boundary}
For every admissible source,
\begin{equation}
 \inf\sigma(H_F|_{\mathcal H_\nu})\downarrow E_\nu
 \quad\hbox{as }F\to\infty.
 \label{flux4-galerkin}
\end{equation}
For a shift $s$ define its outgoing boundary projector by
\[
 B_{s,F}=P_F\,{\mathbf1}_{\{\exists\ell:
                         |E_\ell+s_\ell|>F\}}.
\]
Then, with $Q_F=1-P_F$,
\begin{align}
 Q_FHP_F&=\sum_s U^sA_s B_{s,F},\label{flux4-boundary-identity}\\
 \|Q_FHP_F\psi\|^2
 &\le\mathcal J\sum_s\|A_s\|\|B_{s,F}\psi\|^2
 \le\mathcal J j_{\rm loc}
       \sum_\ell\langle\psi,
          P_F{\bf1}_{\{|E_\ell|=F\}}\psi\rangle
 \quad(\psi=P_F\psi).\label{flux4-boundary-bound}
\end{align}
Nevertheless, already on the source-free physical sector,
\begin{equation}
 \|Q_FHP_F\|\ge\frac1{2g^2a}
 \quad\hbox{for every integer }F\ge0.
 \label{flux4-no-boundary-norm-decay}
\end{equation}
\end{lemma}
\begin{proof}
Finite electric Fourier sums in the Gauss sector are a form core:
$P_F$ commutes with $H_E$ and the Gauss projections, and truncation
converges in the $H_E^{1/2}$ norm.  The remaining terms of $H$ are
bounded.  Rayleigh--Ritz and this form-core approximation prove
\eqref{flux4-galerkin}, including eventual nonemptiness.

The diagonal terms have no outgoing component.  For every shift,
$Q_FU^sA_sP_F=U^sA_sB_{s,F}$, proving the identity.  The triangle
inequality followed by weighted Cauchy--Schwarz proves its first
bound.  Since shifts have coordinates $0,\pm1$, an outgoing transition
must start at $|E_\ell|=F$ for an affected link.  Summing these
indicators and the local coefficient bounds proves the second bound.

For the lower bound choose the electric basis state $e=Fs_p$ and
fill the onsite negative-$\beta$ sea.  The state is normalized,
belongs to $\operatorname{ran}P_F$, and is physical with zero source,
since $\operatorname{div}s_p=0$ and $q_x=0$.  Its outgoing component
at electric configuration $(F+1)s_p$ is the magnetic coefficient
$-b/2$ times that same fermion vector.  No single-link fermion hop
has this shift; any repeated plaquette coefficient has the same sign.
This component alone proves \eqref{flux4-no-boundary-norm-decay}.
\end{proof}

Compression can also create commutators.  For bounded operators $A,B$,
\begin{equation}
 [P_FAP_F,P_FBP_F]
 =P_F[A,B]P_F-P_FAQ_FBP_F+P_FBQ_FAP_F.
 \label{flux4-compressed-commutator}
\end{equation}
Thus even commuting link multipliers can acquire boundary
commutators after compression.  Dropping them would not be a
controlled Schrieffer--Wolff calculation.

\begin{lemma}[A sufficient, volume-dependent flux Schur hypothesis]
\label{flux4-schur}
Fix a Gauss sector, a real spectral parameter $z$, and a nonempty
$P_F$ block.  Write $A=P_FHP_F$, $B=Q_FHP_F$, and let $C$ be
the self-adjoint compression of $H$ to $Q_F$.  A sufficient condition
for $C-z\ge dI>0$ is
\begin{equation}
 d=c(F+1)^2-2Nh_*-z>0.
 \label{flux4-flux-gap-sufficient}
\end{equation}
Under this condition the exact effective operator at $z$ is
\begin{equation}
 A-z-B^*(C-z)^{-1}B,
 \qquad
 0\le B^*(C-z)^{-1}B\le d^{-1}B^*B.
 \label{flux4-flux-schur}
\end{equation}
The original $H-z$ is nonnegative if and only if this effective
operator is nonnegative.  Their null spaces correspond by
$u\mapsto u-(C-z)^{-1}Bu$.
\end{lemma}
\begin{proof}
On $Q_F$, $H_E\ge c(F+1)^2$, and the bounded perturbation gives
$C\ge c(F+1)^2-2Nh_*$.  For $u\in P_F$ and a form-domain
$v\in Q_F$, complete the quadratic form:
\[
 \langle u+v,(H-z)(u+v)\rangle
 =\langle u,[A-z-B^*(C-z)^{-1}B]u\rangle
 +\|(C-z)^{1/2}[v+(C-z)^{-1}Bu]\|^2.
\]
The resolvent bounds and the assertions follow.  All terms preserve
the chosen Gauss sector.
\end{proof}

For $z\le E_0+\Delta$, an explicit sufficient version of
\eqref{flux4-flux-gap-sufficient} is
$c(F+1)^2>\Delta+N[6(1+r)/a+3/(g^2a)]$.
This sufficient condition is not claimed necessary.  It does show
exactly where the available argument retains a volume dependence.

\subsection{What a bounded flux actually proves about pair denominators}

Let $N_{\rm exc}=N_e+N_p$ count particles and holes relative to the
onsite negative-$\beta$ sea.  The diagonal mass-plus-electric
reference is
\begin{equation}
 D=H_E+\mu N_{\rm exc}-2N\mu.
 \label{flux4-diagonal-reference}
\end{equation}
It omits the magnetic plaquette term and the pair-preserving Wilson
hopping.  These terms have not been absorbed into a claimed
diagonal denominator.

\begin{lemma}[The elementary pair denominator on a finite flux cube]
\label{flux4-pair-denominator}
Let $A=tU_\ell^\sigma B_+$ be a gauge-invariant elementary hopping
term that creates one electron--hole pair, where $\sigma=\pm1$,
$B_+$ is a fermion matrix with $[N_{\rm exc},B_+]=2B_+$, and
$\|B_+\|\le1$.  On the electric core,
\begin{equation}
 [D,A]=A\,[2\mu+c(2\sigma E_\ell+1)].
 \label{flux4-exact-denominator}
\end{equation}
For integer $F\ge1$, every nonzero matrix element of
$A_F=P_FAP_F$ creates an energy difference for $D$ at least
\begin{equation}
 \delta_F=2\mu-c(2F-1).
 \label{flux4-internal-denominator}
\end{equation}
In particular $2\mu-c(2F+1)>0$ is a conservative sufficient
condition.  If $\delta_F>0$, there is an anti-Hermitian,
Gauss-preserving operator $S_F$ on $\operatorname{ran}P_F$ with
\begin{equation}
 [S_F,D_F]=-(A_F+A_F^*),\qquad
 \|S_F\|\le\frac{2|t|}{\delta_F},
 \quad D_F=P_FDP_F.
 \label{flux4-first-generator}
\end{equation}
\end{lemma}
\begin{proof}
Since $[E_\ell^2,U_\ell^\sigma]
=U_\ell^\sigma(2\sigma E_\ell+1)$, the first identity follows
by adding the two-excitation mass increment.  For an allowed initial
coordinate $e_\ell$ with $|e_\ell|\le F$, the smallest possible
$2\sigma e_\ell+1$ is $1-2F$.  The transition achieving this
value moves toward the interior, so it is also allowed by the final
cutoff.  This proves the asserted lower bound and its precise
initial-flux convention.

On $P_F$ put
$\Delta_\ell(E)=2\mu+c(2\sigma E_\ell+1)$ and define
$T_F=A_F\Delta_\ell(E)^{-1}$, where the inverse is taken only on
the finite cube.  It is bounded by $\delta_F^{-1}$ and commutes
with $D_F$.  Thus $[D_F,T_F]=A_F$.  Taking adjoints shows that
$S_F=T_F-T_F^*$ has the required commutator and norm.  Both the
hopping and its electric denominator commute with Gauss operators.
\end{proof}

For a sum of elementary pair hops the generators may be added, with
the corresponding sum of norm bounds.  This is an exact first
homological equation for $D_F$, not a block diagonalization of $H_F$.
With $g^2=4\pi\alpha$, the exact elementary positivity condition is
$F<\mu a/(2\pi\alpha)+1/2$; at $\alpha=1/137$ its leading
coefficient is about $21.8\,\mu a$. This large local threshold is
consistent with rare field-driven pairs, but it does not provide the
source-relative global flux budget disproved above.
In particular, the electric-dependent denominator need not commute
with the plaquette term.  The pair-preserving terms and the boundary
commutators in \eqref{flux4-compressed-commutator} still enter the
transformed Hamiltonian.

\paragraph{The additional spectral statement required for pair elimination.}
Inside a fixed-charge, fixed-Gauss $P_F$ space, let $P$ project onto
the chosen no-positron sector and $Q=1-P$.  To apply the same exact
Schur identity to pairs at a real energy $z$, one needs
\begin{equation}
 QH_FQ-zQ\ge d_{\rm pair}Q>0.
 \label{flux4-needed-pair-block}
\end{equation}
It gives a correction
$PH_FQ(QH_FQ-zQ)^{-1}QH_FP$ bounded above by
$d_{\rm pair}^{-1}PH_FQH_FP$.  Any expansion of that inverse
requires a further norm or relative-form bound on the difference
between its actual block and the proposed reference block.
Lemma~\ref{flux4-pair-denominator} supplies neither
\eqref{flux4-needed-pair-block} nor that remainder estimate:
positive differences for individual transitions do not order the
spectra of entire blocks, and the no-pair block contains gauge
excitations.  A charge gap controls $N_e-N_p$, whereas neutral
pairs change $N_e+N_p$ without changing charge.

The results of this section therefore provide valid cutoff
localization, explicit boundary control, finite-volume Galerkin
convergence, and a restricted first pair-elimination equation.
The source-only electric budget is false for this actual model.
A macroscopic low-energy comparison with a source-sized error still
requires dressed-vacuum control and a pair-block resolvent estimate;
it does not follow by deleting the volume term in
\eqref{flux4-volume-budget-eq}.
